\chapter{\MakeUppercase{Methodology and Testing}}

\section{Methodology}
\paragraph{} The approach taken for this project involves a systematic methodology encompassing game development, \ac{rl} environment design, model training, system integration, and evaluation.

\subsection{Game Development Methodology}
\paragraph{} The space shooter game was written using Pygame, a free and open-source library for 2D game programming in Python. The game was designed with modularity in mind, separating player logic, enemy movement, asteroid functionality, and rendering into distinct components within \texttt{game\_3d.py}. An iterative development approach was used, with initial prototypes based around core mechanics like movement and shooting implemented first, followed by enemy \ac{ai}, asteroid spawning, and finally \ac{dda} integration.

\subsection{Reinforcement Learning Methodology}
\paragraph{} The \ac{rl} framework is based on the \ac{ppo} algorithm, primarily chosen for its stable training behavior and sample efficiency. The methodology consists of the following steps:

\begin{enumerate}
	\item \textbf{Environment Definition:} The game environment was wrapped using \texttt{env\_wrapper\_3d.py} to adapt to the Gymnasium interface. The observation space was set as a seven-dimensional vector of player performance indicators.
	
	\item \textbf{State Representation:} Each state $s_t$ at time step $t$ is represented as:
	\begin{equation}
		s_t = [\text{accuracy\_ratio}, \frac{\text{kill\_rate}}{2.0}, \frac{\text{dmg\_rate}}{30.0}, \text{health\_ratio}, \text{shield\_ratio}, \text{wave\_norm}, \text{enemies\_norm}]_t
	\end{equation}
	All metrics are clipped and normalized to the range $[0, 1]$ to facilitate stable training.
	
	\item \textbf{Action Space:} The action space is discrete with five possible actions $a_t \in \{0, 1, 2, 3, 4\}$:
	\begin{itemize}
		\item Action 0: Keep the current level of difficulty the same (no change)
		\item Actions 1, 3: Increase difficulty (enemy count $+1$, enemy speed $+5.0$)
		\item Actions 2, 4: Reduce difficulty (enemy count $-1$, enemy speed $-5.0$)
	\end{itemize}
	
	\item \textbf{Reward Design:} The reward function encourages the agent to maintain player metrics within a target range that reflects the flow state. The reward $r_t$ at time step $t$ is calculated according to the deviation of the player's performance from an ideal engagement zone.
	
	\item \textbf{PPO Training:} The \ac{ppo} algorithm optimizes the policy $\pi_\theta$ by maximizing the clipped surrogate objective:
	\begin{equation}
		L^{CLIP}(\theta) = \hat{\mathbb{E}}_t \left[ \min \left( r_t(\theta) \hat{A}_t, \text{clip}(r_t(\theta), 1-\epsilon, 1+\epsilon) \hat{A}_t \right) \right]
	\end{equation}
	where $r_t(\theta) = \frac{\pi_\theta(a_t|s_t)}{\pi_{\theta_{old}}(a_t|s_t)}$ is the probability ratio, $\hat{A}_t$ is the estimated advantage, and $\epsilon$ is the clipping parameter.
	
	\item \textbf{Model Architecture:} The MLPPolicy3D architecture uses a shared \ac{mlp} backbone with three hidden layers (256 $\rightarrow$ 256 $\rightarrow$ 128 neurons) using Tanh activations. The backbone feeds into an actor head (outputting logits over 5 actions) and a critic head (outputting a scalar value estimate). The network takes the seven-dimensional observation as input.
	
	\item \textbf{Temporal Extension:} The module \texttt{train\_dda\_3d\_lstm.py} extends the base model by adding \ac{lstm} layers to detect temporal trends in player behavior, allowing the agent to take into account past player performances when making difficulty adjustments.
\end{enumerate}

\subsection{Hyperparameter Selection Rationale}
\paragraph{} The stability and convergence speed of the \ac{ppo} algorithm are highly sensitive to hyperparameter configurations. During the training phase, rigorous ablation studies were conducted to converge upon the optimal parameter set for this specific space shooter \ac{dda} task.
\paragraph{} The learning rate ($\alpha$) for both the actor and critic networks was set to a conservatively low $3 \times 10^{-4}$ using the Adam optimizer. While higher learning rates accelerate initial convergence, in a continuous observation space like ours, they frequently led to policy oscillation where the agent would rapidly alternate between maximizing and minimizing difficulty instead of finding a stable equilibrium. The discount factor ($\gamma$) was set to $0.99$, prioritizing long-term rewards over immediate gratification. This ensures the agent does not violently drop the difficulty the moment a player takes minor damage, but rather evaluates the player's health trajectory over a longer horizon.
\paragraph{} The defining feature of \ac{ppo}, the clipping parameter ($\epsilon$), was empirically tuned to $0.20$. This restriction ensures that the newly updated policy does not deviate by more than 20\% from the previous policy during any single update step. Given that dramatic difficulty spikes are notoriously frustrating for players, capping the policy shift was crucial in forcing the agent to learn gradual, smooth adjustments. Additionally, an entropy coefficient ($\beta$) of $0.01$ was applied to the loss function. This intentionally injects a minor degree of randomness into the agent's action selection during early training epochs, effectively preventing premature convergence on suboptimal deterministic policies (e.g., an agent that exclusively always selects "do nothing" to avoid negative rewards).

\subsection{Reward Function Shaping}
\paragraph{} In Reinforcement Learning, the agent is entirely blind to the developer's true intent; it acts purely to maximize the scalar signal provided by the reward function. Therefore, "Reward Shaping" is arguably the most critical methodology in this project. A naive reward function might simply penalize the agent every time the player dies. However, this creates a "sparse reward" problem where the agent receives zero feedback for 99\% of the game, making gradient descent nearly impossible.
\paragraph{} To solve this, a dense, continuous reward function was engineered. The reward $r_t$ is calculated by computing the deviation of the player's current performance $P$ against a predefined optimal target $T$ (the theoretical "flow zone"). The fundamental equation is modeled conceptually as: $r_t = 1.0 - \lambda (P_t - T)^2$, where $\lambda$ is a dynamic scaling penalty factor. 
\paragraph{} If the player's accuracy and survival metrics perfectly align with the target bounds, the agent receives a maximum reward of $+1$. If the player is struggling severely (metrics falling far below the target) or coasting effortlessly (metrics far exceeding the target), the squared deviation results in a negative penalty. Crucially, the reward function was rigorously crafted to prevent "reward hacking." Early training iterations revealed that the agent would occasionally drop the difficulty to absolute zero, allowing the player to easily achieve perfect stats, thus artificially inflating the agent's reward without actually providing engaging gameplay. To counteract this loophole, a secondary penalty term was introduced that directly penalizes the agent mathematically for maintaining ultra-low difficulty configurations for extended, unbroken periods of time $t_{idle}$.

\subsection{Integration Methodology}
\paragraph{} Once the model has been trained, the best-performing version is saved as \texttt{dda\_ppo\_3d.pth}. During live gameplay in \texttt{game\_3d.py}, the model is loaded at initialization. Every few seconds, the current player statistics are gathered, input to the model, and the generated difficulty adjustments are propagated to the game parameters. This inference loop operates asynchronously so as not to affect the game's frame rate.

\section{Testing}
\paragraph{} The testing strategy for this project encompasses unit testing, integration testing, and playtesting.

\subsection{Unit Testing}
\paragraph{} Individual components were tested in isolation:

\begin{itemize}
	\item \textbf{Game Mechanics:} Player movement, shooting, collision detection, and enemy behavior were confirmed through manual and scripted tests.
	\item \textbf{Environment Wrapper:} The observation and action spaces of \texttt{env\_wrapper\_3d.py} were tested to validate correct normalization and bounds.
	\item \textbf{Model Inference:} The trained \ac{ppo} model was run to make sure it produces valid actions within the expected range when given sample observations.
	\item \textbf{API Endpoints:} FastAPI endpoints were tested with automated HTTP requests to verify that data is being logged and retrieved correctly.
\end{itemize}

\subsection{Integration Testing}
\paragraph{} End-to-end integration tests verified that:

\begin{itemize}
	\item The game successfully loads the \ac{ppo} model and performs inference during gameplay.
	\item Difficulty adjustments are applied naturally and do not disrupt the pace of gameplay.
	\item Gameplay data is properly being sent to the FastAPI backend and displayed on the React dashboard.
	\item The system handles edge cases such as model file not found, extreme player performance values, and rapid state transitions.
\end{itemize}

\subsection{Playtesting}
\paragraph{} Numerous playtesting sessions were held with players of varying skill levels to evaluate the effectiveness of the \ac{dda} system. The following aspects were assessed:

\begin{itemize}
	\item Whether the difficulty scales correctly when the player performs well (difficulty should increase).
	\item Whether the challenge gets easier when the player struggles (more health, slower enemies, longer spawn delays).
	\item Overall player satisfaction and engagement during sessions with \ac{dda} enabled versus sessions with a constant level of difficulty.
	\item Stability of the \ac{dda} system — no oscillatory behavior or drastic difficulty swings.
\end{itemize}

\subsection{Extended Testing Strategies}
\paragraph{} Beyond standard unit and integration testing, the system underwent rigorous boundary and stress testing to validate non-functional constraints under extreme conditions.
\paragraph{} \textbf{Stress Testing the API:} The FastAPI telemetry ingestion pipeline was subjected to artificial load testing. Synthetic Python scripts were deployed to bombard the endpoint with 500 concurrent telemetry payloads per second to simulate severe network congestion. The asynchronous ASGI architecture successfully queued all incoming requests without dropping packets or stalling the primary Pygame client loop, validating NFR-01 and NFR-02.
\paragraph{} \textbf{Boundary Condition Testing:} The \ac{dda} mathematical bounds were aggressively tested by simulating robotic "perfect" players and "AFK" (Away From Keyboard) players. For perfect players (100\% accuracy, 0 damage taken), the network successfully scaled the difficulty to the maximum hardcoded ceiling within 45 seconds without throwing buffer overflow or out-of-bounds array exceptions. For AFK players, the system rapidly dropped enemies to the minimum spawn rate, validating the mathematical safety nets of the normalization functions.
\paragraph{} \textbf{User Acceptance Testing (UAT) Rubrics:} For human playtesting, subjective analysis was quantified using standardized UAT rubrics loosely based on the Game Experience Questionnaire (GEQ). Players were asked to individually rate their experiences on a 5-point Likert scale across dimensions including Competence, Immersion, Flow, Tension, and Challenge. These qualitative survey scores were directly cross-referenced with the quantitative telemetry logged by the FastAPI backend to definitively prove the real-world efficacy of the \ac{rl} agent from a psychological perspective.
